% ==========================================================================
%  Chapter 72 -- Mixed-Model Reinforced Sub-Model with Discrete Rebars
% ==========================================================================

\chapter{Mixed-Model Reinforced Sub-Model with Discrete Rebars}
\label{ch:mixed-model-rebar}


% =========================================================================
\section{Introduction and Motivation}
\label{sec:mixed-intro}

The persistent nonlinear sub-model evolver introduced in
Chapter~\ref{ch:nl-sub-model-evolver} uses a homogeneous three-dimensional
continuum model (all Hex27 elements with the Ko--Bathe concrete
material) to track damage evolution at critical locations.  While this
captures concrete cracking and damage realistically, it neglects the
contribution of reinforcing steel, which dominates the post-cracking
response of reinforced-concrete (RC) members.

This chapter extends the sub-model framework to support \emph{mixed
models} that embed discrete one-dimensional rebar elements (two-node
trusses) within the three-dimensional hexahedral mesh.  The resulting
model has two element types sharing nodes through the DMPlex sieve
(perfect-bond assumption), a mixed finite-element configuration that
cannot be handled by the single-type \texttt{SingleElementPolicy}.  The
type-erased \texttt{FEM\_Element} wrapper and the
\texttt{MultiElementPolicy} provide the mechanism to store and iterate
over heterogeneous element types in a single \texttt{Model<>}.


% =========================================================================
\section{Mathematical Formulation}
\label{sec:mixed-math}

% -------------------------------------------------------------------------
\subsection{Mixed Displacement Field and Perfect Bond}

Let $\Omega_c \subset \mathbb{R}^3$ denote the concrete domain
discretised by hexahedral elements, and let $\Gamma_s \subset \Omega_c$
denote the set of one-dimensional rebar segments embedded within the
concrete mesh.  Under the perfect-bond assumption, the rebar
displacement at any point along the bar coincides with the continuum
displacement interpolated at the shared nodes:
\begin{equation}
  \label{eq:perfect-bond}
  \mathbf{u}^s(\mathbf{x}) = \mathbf{u}^c(\mathbf{x}), \quad
  \forall\; \mathbf{x} \in \Gamma_s \cap \partial\Omega_c^{\text{nodes}},
\end{equation}
where the superscripts $c$ and $s$ refer to concrete and steel,
respectively.  Since truss nodes coincide with hexahedral nodes the
shared DOFs enter both element stiffness contributions additively.

% -------------------------------------------------------------------------
\subsection{Continuum Element Contribution (Concrete)}

Each hexahedral element $e$ contributes the standard 3D solid tangent
stiffness and internal-force vector:
\begin{equation}
  \label{eq:hex-Kt}
  \mathbf{K}_t^{(e)} = \int_{\Omega^{(e)}}
    \mathbf{B}_c^T \, \mathbf{D}_t^{(e)} \, \mathbf{B}_c \; \mathrm{d}\Omega,
  \qquad
  \mathbf{f}_{\text{int}}^{(e)} = \int_{\Omega^{(e)}}
    \mathbf{B}_c^T \, \boldsymbol{\sigma}^{(e)} \; \mathrm{d}\Omega,
\end{equation}
where $\mathbf{B}_c$ is the standard strain--displacement matrix,
$\mathbf{D}_t$ the consistent tangent modulus from the Ko--Bathe
material, and $\boldsymbol{\sigma}$ the Cauchy stress vector.  The
integrals are evaluated via Gauss--Legendre quadrature at
$3 \times 3 \times 3 = 27$ points for Hex27 elements.

% -------------------------------------------------------------------------
\subsection{Truss Element Contribution (Rebar)}

Each rebar segment is modelled as a two-node bar element in 3D space
with $2 \times 3 = 6$ DOFs.  The element strain is purely axial:
\begin{equation}
  \label{eq:truss-eps}
  \varepsilon = \mathbf{B}_s \, \mathbf{u}_e,
  \qquad
  \mathbf{B}_s = \frac{1}{L}
  \begin{bmatrix} -l & -m & -n & l & m & n \end{bmatrix},
\end{equation}
where $(l, m, n)$ are the direction cosines of the bar axis and $L$ is
the element length.  The tangent stiffness and internal force are:
\begin{equation}
  \label{eq:truss-Kt}
  \mathbf{K}_t^{(s)} = A_s \, L \, E_t(\varepsilon)\;
  \mathbf{B}_s^T \, \mathbf{B}_s,
  \qquad
  \mathbf{f}_{\text{int}}^{(s)} = A_s \, L \, \sigma(\varepsilon)\;
  \mathbf{B}_s^T,
\end{equation}
where $A_s$ is the bar cross-sectional area and
$E_t = \partial\sigma / \partial\varepsilon$ is
the uniaxial tangent modulus from the Menegotto--Pinto steel model.

One Gauss point suffices for exact integration over a 2-node linear
element with constant $\mathbf{B}_s$.

% -------------------------------------------------------------------------
\subsection{Assembled Mixed System}

The global stiffness matrix and internal-force vector are obtained by
the standard direct-stiffness assembly:
\begin{equation}
  \label{eq:mixed-assembly}
  \mathbf{K}_t = \sum_{e \in \mathcal{E}_c} \mathbf{A}_e^T
    \mathbf{K}_t^{(e)} \mathbf{A}_e
  + \sum_{s \in \mathcal{E}_s} \mathbf{A}_s^T
    \mathbf{K}_t^{(s)} \mathbf{A}_s,
\end{equation}
where $\mathbf{A}_e$ and $\mathbf{A}_s$ are the Boolean scatter
matrices that map element to global DOFs.  Because rebar nodes are
shared with hex nodes, the PETSc \texttt{MatSetValuesLocal} assembly
with \texttt{ADD\_VALUES} automatically couples both contributions at
shared DOFs.

The equilibrium equation at each Newton iteration is therefore:
\begin{equation}
  \label{eq:mixed-Newton}
  \mathbf{K}_t\;\delta\mathbf{u} = -\mathbf{R},
  \qquad
  \mathbf{R} = \mathbf{f}_{\text{int}}^c + \mathbf{f}_{\text{int}}^s
               - \mathbf{f}_{\text{ext}},
\end{equation}
where $\mathbf{f}_{\text{int}}^c$ and $\mathbf{f}_{\text{int}}^s$
come from concrete and rebar elements respectively.


% =========================================================================
\section{Type-Erased Element Access: \texttt{GaussPointSnapshot}}
\label{sec:mixed-gauss-snapshot}

With \texttt{MultiElementPolicy}, the model stores elements as
\texttt{std::vector<FEM\_Element>}, where \texttt{FEM\_Element} is a
type-erased wrapper using the External Polymorphism pattern.  The
heterogeneous element list means we cannot call typed methods like
\texttt{material\_points()} directly.

A new virtual method \texttt{gauss\_point\_snapshots(Vec u\_local)} was
added to the \texttt{FEM\_Element::Concept} base class.  The
\texttt{Model<T>} bridge template implements this method using
\texttt{if constexpr}:

\begin{itemize}
  \item For \texttt{ContinuumElement}: iterates \texttt{material\_points()},
    extracts position from geometry integration points, stress/strain from
    the material state, and crack data from
    \texttt{internal\_field\_snapshot()}.
  \item For \texttt{TrussElement}: returns an empty vector (steel has no
    crack model; its contribution is purely structural).
\end{itemize}

The returned \texttt{GaussPointSnapshot} struct aggregates:
\begin{equation*}
  \texttt{GaussPointSnapshot} = \bigl\{
    \mathbf{x}, \;
    \boldsymbol{\sigma}, \;
    \boldsymbol{\varepsilon}, \;
    d, \;
    n_{\text{cracks}}, \;
    \{\mathbf{n}_i\}_{i=1}^3, \;
    \{w_i\}_{i=1}^3, \;
    \{c_i\}_{i=1}^3
  \bigr\},
\end{equation*}
where $d$ is the scalar damage, $\mathbf{n}_i$ are crack normals,
$w_i$ are crack openings, and $c_i$ are open/closed flags.  This
struct provides a uniform interface for both crack-data collection
and volume-averaged stress/strain extraction.


% =========================================================================
\section{Reinforced Prismatic Domain Construction}
\label{sec:mixed-domain}

The function \texttt{make\_reinforced\_prismatic\_domain()} generates
a combined hex+line mesh:

\begin{enumerate}
  \item All nodes are created on the structured prismatic grid
    (same as the homogeneous case).
  \item Hexahedral element geometries are created with the specified
    element order (Hex8, Hex20, or Hex27), each assigned to the
    \texttt{"Solid"} physical group.
  \item For each rebar bar specified by grid indices $(i_x, i_y)$
    and cross-sectional area $A_s$, the builder creates $n_z$ 2-node
    line elements connecting consecutive z-level nodes at
    $(\mathtt{step} \cdot i_x, \; \mathtt{step} \cdot i_y, \;
    \mathtt{step} \cdot k)$ for $k = 0, \ldots, n_z - 1$.
    These elements are assigned to the \texttt{"Rebar"} physical group.
  \item A \texttt{RebarElementRange} is returned, providing the index
    range $[\mathtt{first}, \mathtt{last})$ of rebar element
    geometries within the domain's element list.
\end{enumerate}

The coordinator maps physical rebar positions $(y, z)$ from the
\texttt{SubModelSpec} to grid indices via nearest-node snapping:
\begin{equation}
  \label{eq:rebar-snap}
  i_x = \operatorname{round}\!\left(
    \frac{y + w/2}{\Delta x}\right), \qquad
  i_y = \operatorname{round}\!\left(
    \frac{z + h/2}{\Delta y}\right),
\end{equation}
where $w$ and $h$ are the section width and height, and
$\Delta x = w / n_x$, $\Delta y = h / n_y$ are the element spacings.
The indices are clamped to $[0, n_x]$ and $[0, n_y]$ respectively.


% =========================================================================
\section{Model Construction in the Evolver}
\label{sec:mixed-evolver}

The \texttt{NonlinearSubModelEvolver} was refactored from a
\texttt{HomogModel} (typed, \texttt{SingleElementPolicy}) to a
\texttt{MixedModel} (\texttt{MultiElementPolicy}) which always stores
elements as \texttt{FEM\_Element}.  The \texttt{first\_solve()} method
builds the element list explicitly:

\begin{enumerate}
  \item \textbf{Concrete elements}: For each hex geometry
    (indices $[0, \mathtt{rebar\_first})$), create a
    \texttt{ContinuumElement} with the Ko--Bathe concrete material and
    wrap in \texttt{FEM\_Element}.
  \item \textbf{Rebar elements} (optional): For each line geometry
    (indices $[\mathtt{rebar\_first}, \mathtt{rebar\_last})$), create a
    \texttt{TrussElement<3>} with the Menegotto--Pinto steel material,
    the bar's cross-sectional area, and wrap in \texttt{FEM\_Element}.
  \item Pass the assembled \texttt{std::vector<FEM\_Element>} to the
    Model constructor~2 (pre-built elements), then proceed with
    constraint application and PETSc SNES setup as before.
\end{enumerate}

The SNES callbacks (\texttt{FormResidual}, \texttt{FormJacobian})
iterate \texttt{model\_->elements()} and call
\texttt{compute\_internal\_forces}, \texttt{inject\_tangent\_stiffness},
\texttt{commit\_material\_state}, and \texttt{revert\_material\_state}
--- all of which are forwarded through the \texttt{FEM\_Element}
type-erasure to the underlying concrete or truss element.

\subsection{Post-Processing Adaptation}

The \texttt{extract\_results()} and \texttt{collect\_crack\_data()}
methods were rewritten to use \texttt{gauss\_point\_snapshots()}
instead of the typed \texttt{material\_points()} access:
\begin{itemize}
  \item \texttt{extract\_results()}: Accumulates volume-averaged
    $\boldsymbol{\sigma}$ and $\boldsymbol{\varepsilon}$ from
    snapshots, computes effective $E$ and $G$, and evaluates the
    von~Mises peak.
  \item \texttt{collect\_crack\_data()}: Iterates snapshots with
    $n_{\text{cracks}} > 0$ to build \texttt{CrackRecord} entries for
    VTK crack-plane export.
\end{itemize}
Truss elements return empty snapshot vectors, so they contribute
neither to crack statistics nor to the stress averaging (by design,
rebar stress is a uniaxial contribution embedded in the continuum
average).


% =========================================================================
\section{Coordinator Integration}
\label{sec:mixed-coord}

The \texttt{MultiscaleCoordinator::build\_sub\_models()} method was
extended to detect \texttt{spec.has\_rebar()} and dispatch accordingly:
\begin{itemize}
  \item \textbf{Without rebar}: calls \texttt{make\_prismatic\_domain()}
    as before, producing a homogeneous hex mesh.
  \item \textbf{With rebar}: converts the physical rebar coordinates to
    grid indices via \eqref{eq:rebar-snap}, builds a
    \texttt{RebarSpec}, and calls
    \texttt{make\_reinforced\_prismatic\_domain()}.
\end{itemize}

The \texttt{MultiscaleSubModel} struct was extended with:
\begin{itemize}
  \item \texttt{rebar\_range}: the element index range for rebar
    geometries within the domain.
  \item \texttt{rebar\_areas}: a vector of bar areas (one per bar,
    repeated for each z-level segment internally).
  \item \texttt{has\_rebar()}: convenience query.
\end{itemize}

The evolver's \texttt{set\_rebar\_material(E, fy, b)} method allows
specifying the Menegotto--Pinto parameters (elastic modulus, yield
stress, hardening ratio) before the first solve.


% =========================================================================
\section{Summary}
\label{sec:mixed-summary}

The mixed-model extension enables realistic simulation of
reinforced-concrete sub-models at the microscale:
\begin{enumerate}
  \item Concrete (Hex27 + Ko--Bathe): captures cracking, softening,
    and damage accumulation.
  \item Steel (TrussElement<3> + Menegotto--Pinto): captures yielding,
    Bauschinger effect, and cyclic hardening.
  \item Perfect bond: shared nodes enforce displacement compatibility
    through the assembled stiffness.
  \item Type erasure: \texttt{FEM\_Element} and
    \texttt{MultiElementPolicy} allow both element types in a single
    \texttt{Model<>} with unified SNES callbacks.
\end{enumerate}

The implementation preserves full backward compatibility: when
\texttt{SubModelSpec} has no rebar bars, the evolver creates the
same homogeneous concrete model as before.
